\mysec{A medida de Lebesgue em $\mathbb{R}$}

\noindent
Nosso objetivo será construir uma medida $m$ em $\mathbb{R}$ que satisfaz as seguintes propriedades:
\begin{itemize}
    \item[i)] $m(I) = b-a$ (comprimento de um intervalo $I = (a,b), (a,b], [a,b)$ ou $[a,b]$) e $m(I) = +\infty$ se $I$ for um intervalo ilimitado.
    \item[ii)] $m(E+y) = m(E)$ (a medida é invariante para translações).
    \item[iii)] $m(\emptyset) = 0$; $m(E) \ge 0$ e $m\left(\bigcup E_k\right) = \sum m(E_k)$ (é de fato uma medida).
\end{itemize}

Infelizmente tal medida não existe ao considerar $\mathcal{P}(\mathbb{R})$ como domínio. Primeiro construiremos uma $\sigma$-álgebra em que (i), (ii) e (iii) serão satisfeitas.

\mysubsec{Medida Exterior de Lebesgue}

Para um intervalo $I$ de $\mathbb{R}$, definimos seu comprimento $l(I) = b-a$, se $I = (a,b), (a,b], [a,b), [a,b]$, e $l(I) = +\infty$ se $I$ não for limitado.

\begin{definition*}{}
Para um conjunto $A \in \mathcal{P}(\mathbb{R})$, definimos sua medida exterior como
$$ m^*(A) = \inf \left\{ \sum_{k \in \mathbb{N}} l(I_k) ; A \subseteq \bigcup_{k \in \mathbb{N}} I_k \right\} $$
(Podemos restringir para que todos os $I_k$ sejam abertos).
\end{definition*}

Segue imediatamente da definição que $m^*(A) \ge 0, \forall A \in \mathcal{P}(\mathbb{R})$, $m^*(\emptyset) = 0$ e que se $A \subseteq B$, então $m^*(A) \le m^*(B)$, uma vez que toda cobertura por intervalos de $B$ também cobrirá $A$.

\textbf{Exemplo:} Todo conjunto enumerável possui medida exterior zero. De fato, se $C = \{c_k\}_{k \in \mathbb{N}}$, dado $\epsilon > 0$ defina
$$ I_k := \left(c_k - \frac{\epsilon}{2^{k+1}}, c_k + \frac{\epsilon}{2^{k+1}}\right). $$
Assim, $C \subseteq \bigcup_{k \in \mathbb{N}} I_k$ e $\sum_{k \in \mathbb{N}} l(I_k) = \sum_{k \in \mathbb{N}} \frac{\epsilon}{2^k} = \epsilon \implies m^*(C) \le \epsilon$. Como vale para todo $\epsilon > 0$, obtemos $m^*(C) = 0$.

\begin{prop*}{1}
  Para todo intervalo $I$ de $\mathbb{R}$, $\extmu{I} = l(I)$.
\end{prop*}

\begin{proof}[\textbf{Dem:}]
Comecemos com o caso $I = [a,b]$. Dado $\epsilon > 0$, temos que $[a,b] \subset (a-\epsilon, b+\epsilon)$. Assim, $m^*([a,b]) \le b-a+2\epsilon$. Como vale para todo $\epsilon > 0$, temos que $m^*([a,b]) \le b-a$.

Vejamos agora que $m^*([a,b]) \ge b-a$. Para isso, mostraremos que $b-a$ é uma cota inferior, ou seja, $\sum_{k \in \mathbb{N}} l(I_k) \ge b-a$, em que $\{I_k\}_{k \in \mathbb{N}}$ é uma cobertura enumerável por intervalos abertos de $[a,b]$, com $[a,b] \subseteq \bigcup_{k \in \mathbb{N}} I_k$.

Como $[a,b]$ é compacto, existem finitos intervalos destes que cobrem $[a,b]$. Denotaremos por $\{I_k\}_{k=1}^n$. 
Como $a \in [a,b]$, $a$ pertence a algum desses intervalos, que denotaremos por $(a_1, b_1)$, com $a_1 < a < b_1$.

Se $b_1 > b$, então $\sum_{k=1}^n l(I_k) \ge b_1 - a_1 > b - a$.
Se $b_1 \le b$, então $b_1 \in [a,b]$ e, portanto, existe um intervalo na coleção que contém $b_1$. Chamaremos de $(a_2, b_2)$, com $a_2 < b_1 < b_2$. Então
$$ \sum_{k=1}^n l(I_k) \ge (b_1 - a_1) + (b_2 - a_2) = b_2 + (b_1 - a_2) - a_1 > b_2 - a_1. $$
Se $b_2 \le b$, podemos continuar esse processo até encontrar $b_N > b$. Uma vez que temos um número finito de intervalos $\{I_k\}_{k=1}^n$, o processo termina.

Assim, obtemos uma subcoleção $\{(a_j, b_j)\}_{j=1}^N$ de $\{I_k\}$ tal que $a_1 < a$, $a_j < b_{j-1} < b_j$ para $j=2, \dots, N$ e $b_N > b$.
Portanto,
\begin{align*}
    \sum_{k \in \mathbb{N}} l(I_k) &\ge \sum_{k=1}^n l(I_k) \ge \sum_{j=1}^N l(a_j, b_j) = (b_N - a_N) + \dots + (b_1 - a_1) \\
    &= b_N + (b_{N-1} - a_N) + (b_{N-2} - a_{N-1}) + \dots + (b_1 - a_2) - a_1 \\
    &> b_N - a_1 > b - a.
\end{align*}
Com isso, $b-a \le m^*([a,b])$.

Se $I$ é um intervalo limitado, dado $\epsilon > 0$, existem intervalos fechados $J_1$ e abertos $J_2$ tais que $J_1 \subseteq I \subseteq J_2$ com $l(I) - \epsilon < l(J_1) \le l(J_2) < l(I) + \epsilon$.
Assim,
$$ l(I) - \epsilon < l(J_1) = m^*(J_1) \le m^*(I) \le m^*(J_2) = l(J_2) < l(I) + \epsilon. $$
Como vale para todo $\epsilon > 0$, obtemos $m^*(I) = l(I)$.

Por fim, se $I$ é ilimitado, para cada $n \in \mathbb{N}$, existe um intervalo $J \subseteq I$ com $l(J) = n$. Assim, $m^*(I) \ge m^*(J) = l(J) = n \implies m^*(I) = +\infty = l(I)$.
\end{proof}

\begin{prop*}{2}
A medida exterior $m^*$ é invariante por translações, isto é, $m^*(A+y) = m^*(A)$, em que $A+y = \{a+y ; a \in A\}$.
\end{prop*}

\begin{proof}[\textbf{Dem:}]
Note que $\{I_k\}_{k \in \mathbb{N}}$ é uma cobertura por intervalos de $A$ se, e somente se, $\{I_k+y\}_{k \in \mathbb{N}}$ é uma cobertura para $A+y$.
Além disso, $l(I_k) = l(I_k+y)$. Portanto $\sum_{k \in \mathbb{N}} l(I_k) = \sum_{k \in \mathbb{N}} l(I_k+y)$. Logo, $m^*(A) = m^*(A+y)$.
\end{proof}

\begin{prop*}{3}
A medida exterior é enumeravelmente subaditiva, isto é, se $\{E_k\}_{k \in \mathbb{N}}$ é qualquer coleção enumerável de conjuntos, então
$$ m^*\left(\bigcup_{k \in \mathbb{N}} E_k\right) \le \sum_{k \in \mathbb{N}} m^*(E_k). $$
\end{prop*}

\begin{proof}[\textbf{Dem:}]
Se algum $E_k$ satisfizer $m^*(E_k) = +\infty$, então a desigualdade segue trivialmente. Suponhamos então que $m^*(E_k) < +\infty$, para todo $k \in \mbb N$. 

Dado $\epsilon > 0$, para cada $k \in \mathbb{N}$, existe uma coleção enumerável de intervalos $\{I_j^k\}_{j \in \mathbb{N}}$ tais que $E_k \subseteq \bigcup_{j \in \mathbb{N}} I_j^k$ com
  $$ \sum_{j \in \mathbb{N}} l(I_j^k) < m^*(E_k) + \frac{\epsilon}{2^k}. $$
Assim, $\{I_j^k\}_{j,k \in \mathbb{N}}$ é uma coleção enumerável de intervalos. Como $\bigcup_{k \in \mathbb{N}} E_k \subseteq \bigcup_{k \in \mathbb{N}} \bigcup_{j \in \mathbb{N}} I_j^k$, temos
$$ m^*\left(\bigcup_{k \in \mathbb{N}} E_k\right) \le \sum_{k \in \mathbb{N}} \sum_{j \in \mathbb{N}} l(I_j^k) < \sum_{k \in \mathbb{N}} \left(m^*(E_k) + \frac{\epsilon}{2^k}\right) = \sum_{k \in \mathbb{N}} m^*(E_k) + \epsilon. $$
Como vale para todo $\epsilon > 0$, concluímos que $m^*\left(\bigcup_{k \in \mathbb{N}} E_k\right) \le \sum_{k \in \mathbb{N}} m^*(E_k)$.

A subaditividade finita segue da proposição anterior ao tomar $E_k = \emptyset$ para $k > n$, obtendo:
$$ m^*\left(\bigcup_{k=1}^n E_k\right) = m^*\left(\bigcup_{k \in \mathbb{N}} E_k\right) \le \sum_{k \in \mathbb{N}} m^*(E_k) = \sum_{k=1}^n m^*(E_k). $$
\end{proof}

\mysubsec{Conjuntos Lebesgue Mensuráveis}

Embora $m^*$ possua todas essas boas propriedades, ela falha em ser aditiva! Veremos nas próximas aulas que existem conjuntos disjuntos $A$ e $B$ tais que $m^*(A \cup B) < m^*(A) + m^*(B)$. Assim, precisaremos restringir o domínio de $m^*$ para conseguirmos a propriedade da aditividade enumerável.

\begin{definition*}{}
Dizemos que $E \in \mathcal{P}(\mathbb{R})$ é (Lebesgue) mensurável se
$$ m^*(A) = m^*(A \cap E) + m^*(A \cap E^c), \quad \forall A \in \mathcal{P}(\mathbb{R}) $$
(Condição de Carathéodory).
\end{definition*}

\noindent
Note que se $A$ e $B$ forem disjuntos e um deles for mensurável (digamos $A$ mensurável), então
$$ m^*(A \cup B) = m^*((A \cup B) \cap A) + m^*((A \cup B) \cap A^c) = m^*(A) + m^*(B). $$
Além disso, note que $A = (A \cap E) \cup (A \cap E^c), \forall A \in \mathcal{P}(\mathbb{R})$, e portanto,
$$ m^*(A) \le m^*(A \cap E) + m^*(A \cap E^c). $$
Assim, para verificar se $E$ é mensurável, basta apenas mostrar que
$$ m^*(A) \ge m^*(A \cap E) + m^*(A \cap E^c), \quad \forall A \in \mathcal{P}(\mathbb{R}). $$
Como a desigualdade é satisfeita para $m^*(A) = +\infty$, basta se restringir para conjuntos com $m^*(A) < +\infty$.

\begin{prop*}{4}
Temos que $\emptyset$ e $\mathbb{R}$ são mensuráveis e se $E$ é mensurável, então $E^c$ é mensurável.
\end{prop*}

\begin{proof}[\textbf{Dem:}]
$$ m^*(A \cap \emptyset) + m^*(A \cap \emptyset^c) = m^*(\emptyset) + m^*(A) = m^*(A), \quad \forall A \in \mathcal{P}(\mathbb{R}). $$
$$ m^*(A \cap \mathbb{R}) + m^*(A \cap \mathbb{R}^c) = m^*(A) + m^*(\emptyset) = m^*(A), \quad \forall A \in \mathcal{P}(\mathbb{R}). $$
Como $(E^c)^c = E$, segue que se $E$ é mensurável, então
$$ m^*(A \cap E^c) + m^*(A \cap (E^c)^c) = m^*(A \cap E^c) + m^*(A \cap E) = m^*(A), \quad \forall A \in \mathcal{P}(\mathbb{R}). $$
\end{proof}

\begin{prop*}{5}
Se $m^*(E) = 0$, então $E$ é mensurável.
\end{prop*}

\begin{proof}[\textbf{Dem:}]
Seja $A \in \mathcal{P}(\mathbb{R})$. Como $A \cap E \subseteq E$, então $m^*(A \cap E) \le m^*(E) = 0 \implies m^*(A \cap E) = 0$.
Além disso, $A \cap E^c \subseteq A \implies m^*(A \cap E^c) \le m^*(A)$.
Assim,
$$ m^*(A \cap E) + m^*(A \cap E^c) \le 0 + m^*(A) = m^*(A). $$
Portanto $E$ é mensurável.
\end{proof}

\begin{prop*}{6}
A união finita de conjuntos mensuráveis é mensurável.
\end{prop*}

\begin{proof}[\textbf{Dem:}]
Inicialmente, provaremos que se $E_1, E_2$ são mensuráveis, então $E_1 \cup E_2$ é mensurável. Seja $A \in \mathcal{P}(\mathbb{R})$. Então
$$ m^*(A) = m^*(A \cap E_1) + m^*(A \cap E_1^c) = m^*(A \cap E_1) + m^*([A \cap E_1^c] \cap E_2) + m^*([A \cap E_1^c] \cap E_2^c) $$
Como $(A \cap E_1^c) \cap E_2^c = A \cap (E_1 \cup E_2)^c$ e $(A \cap E_1) \cup (A \cap E_1^c \cap E_2) = A \cap (E_1 \cup E_2)$, temos pela subaditividade que
$$ m^*(A \cap (E_1 \cup E_2)) + m^*(A \cap (E_1 \cup E_2)^c) \le m^*(A \cap E_1) + m^*(A \cap E_1^c \cap E_2) + m^*(A \cap E_1^c \cap E_2^c) = m^*(A). $$
Portanto $E_1 \cup E_2$ é mensurável. O caso geral segue por indução no número de conjuntos.
\end{proof}

\begin{prop*}{7}
Seja $A \in \mathcal{P}(\mathbb{R})$ e $\{E_k\}_{k=1}^n$ uma coleção finita de conjuntos disjuntos mensuráveis. Então
$$ m^*\left(A \cap \left(\bigcup_{k=1}^n E_k\right)\right) = \sum_{k=1}^n m^*(A \cap E_k). $$
Em particular, $m^*\left(\bigcup_{k=1}^n E_k\right) = \sum_{k=1}^n m^*(E_k)$.
\end{prop*}

\begin{proof}[\textbf{Dem:}]
Provaremos por indução em $n$. O caso $n=1$ é trivial. Suponhamos a validade para $n-1$ e considere $\{E_k\}_{k=1}^n$ conjuntos disjuntos mensuráveis.
Como eles são disjuntos, $A \cap \left(\bigcup_{k=1}^n E_k\right) \cap E_n = A \cap E_n$ e $A \cap \left(\bigcup_{k=1}^n E_k\right) \cap E_n^c = A \cap \left(\bigcup_{k=1}^{n-1} E_k\right)$.
Assim, como $E_n$ é mensurável, obtemos
\begin{align*}
    m^*\left(A \cap \left(\bigcup_{k=1}^n E_k\right)\right) &= m^*\left(A \cap \left(\bigcup_{k=1}^n E_k\right) \cap E_n\right) + m^*\left(A \cap \left(\bigcup_{k=1}^n E_k\right) \cap E_n^c\right) \\
    &= m^*(A \cap E_n) + m^*\left(A \cap \left(\bigcup_{k=1}^{n-1} E_k\right)\right) \\
    &= m^*(A \cap E_n) + \sum_{k=1}^{n-1} m^*(A \cap E_k) = \sum_{k=1}^n m^*(A \cap E_k).
\end{align*}
A última igualdade segue pela hipótese de indução.
\end{proof}

\begin{prop*}{8}
A união enumerável de conjuntos mensuráveis é mensurável.
\end{prop*}

\begin{proof}[\textbf{Dem:}]
Seja $\{B_k\}_{k \in \mathbb{N}}$ uma coleção enumerável de conjuntos mensuráveis. Pela mesma construção do lema da aula passada, existe uma coleção enumerável disjunta $\{E_k\}_{k \in \mathbb{N}}$ mensurável tal que $E_k \subseteq B_k, \forall k \in \mathbb{N}$ e $\bigcup_{k \in \mathbb{N}} B_k = \bigcup_{k \in \mathbb{N}} E_k = E$.

Para cada $n \in \mathbb{N}$, defina $F_n = \bigcup_{k=1}^n E_k$. Assim, cada $F_n$ é mensurável e $F_n \subseteq E \implies E^c \subseteq F_n^c$.
Assim, para $A \in \mathcal{P}(\mathbb{R})$,
$$ m^*(A) = m^*(A \cap F_n) + m^*(A \cap F_n^c) \ge m^*(A \cap F_n) + m^*(A \cap E^c). $$
Pela Proposição 7, $m^*(A \cap F_n) = \sum_{k=1}^n m^*(A \cap E_k)$. Assim, 
$$ m^*(A) \ge \sum_{k=1}^n m^*(A \cap E_k) + m^*(A \cap E^c). $$
Como essa desigualdade vale para todo $n \in \mathbb{N}$, temos que $m^*(A) \ge \sum_{k=1}^{+\infty} m^*(A \cap E_k) + m^*(A \cap E^c)$.
Assim, pela subaditividade enumerável de $m^*$, obtemos
$$ m^*(A \cap E) + m^*(A \cap E^c) \le \sum_{k \in \mathbb{N}} m^*(A \cap E_k) + m^*(A \cap E^c) \le m^*(A). $$
Portanto $E$ é mensurável.
\end{proof}

\bigline
\noindent
\textbf{Conclusão:} A coleção $\mathcal{M} = \{E \in \mathcal{P}(\mathbb{R}) ; E \text{ é Lebesgue mensurável}\}$ é uma $\sigma$-álgebra.
